% main-ieee.tex — Paper 2: Unilateral WiFi CSI as NIST-Validated Entropy Source
% IEEE format for IACR ePrint submission
% Converted from ACM sigconf (main.tex)
% Author: Daniel Mo Houshmand, QDaria Quantum Research

\documentclass[conference]{IEEEtran}

%% ─── Packages ───
\usepackage{cite}
\usepackage{amsmath,amssymb,amsfonts,amsthm}
\usepackage{algorithm}
\usepackage{algorithmic}
\usepackage{booktabs}
\usepackage{siunitx}
\usepackage{subcaption}
\usepackage{xspace}
\usepackage{balance}
\usepackage{tikz}
\usepackage{pgfplots}
\pgfplotsset{compat=1.18}
\usetikzlibrary{arrows.meta,positioning,shapes.geometric,calc,decorations.pathreplacing}
\usepackage{graphicx}
\usepackage{url}
\usepackage{hyperref}
\usepackage{microtype}

%% ─── Theorem-like environments ───
\newtheorem{theorem}{Theorem}
\newtheorem{corollary}[theorem]{Corollary}
\newtheorem{proposition}[theorem]{Proposition}
\newtheorem{lemma}[theorem]{Lemma}
\newtheorem{definition}[theorem]{Definition}
\newtheorem{assumption}{Assumption}
\newtheorem*{remark}{Remark}

%% ─── Custom commands ───
\newcommand{\zipminator}{\textsc{Zipminator}\xspace}
\newcommand{\puek}{\textsc{PUEK}\xspace}
\newcommand{\minentropy}{H_{\infty}}
\newcommand{\shannon}{H}
\newcommand{\bpb}{\text{bits/byte}}
\newcommand{\eg}{e.g.,\xspace}
\newcommand{\ie}{i.e.,\xspace}
\newcommand{\etal}{et~al.\xspace}

\graphicspath{{./figures/}}

\begin{document}

\title{Unilateral WiFi CSI as a NIST-Validated Entropy Source: \\
From Bilateral Key Agreement to Single-Device Randomness}

\author{\IEEEauthorblockN{Daniel Mo Houshmand}
\IEEEauthorblockA{QDaria Quantum Research, Oslo, Norway\\
\texttt{mo@qdaria.com}\\
ORCID: 0009-0008-2270-5454}}

\maketitle

\begin{abstract}
Every prior system that extracts randomness from WiFi Channel State Information (CSI) requires two cooperating endpoints exploiting channel reciprocity for bilateral key agreement.
We present the first system, measurement, and NIST SP~800-90B validation of WiFi CSI as a \emph{unilateral} entropy source: a single device passively measuring ambient CSI to harvest genuine physical randomness with no cooperating partner.
Using the public Gi-z/CSI-Data corpus (TU~Darmstadt Nexmon captures, Broadcom BCM4339), we extract phase least-significant bits from 343~frames across 256~OFDM subcarriers, apply Von~Neumann debiasing, and obtain 2,690~bytes of entropy at a 24.5\% extraction ratio.
The NIST SP~800-90B \texttt{ea\_non\_iid} assessment yields a final min-entropy of \textbf{5.50~bits/byte} (MCV estimator, 99\% confidence), compared to 6.35 for IBM Quantum (ibm\_kingston, 156~qubits) and 6.36 for \texttt{os.urandom}.
We introduce the Physical Unclonable Environment Key (PUEK), which derives location-locked cryptographic keys from the SVD eigenstructure of CSI measurements, with security profiles from $\tau=0.75$ (office) to $\tau=0.98$ (military).
The method is hardware-agnostic, working with any IEEE~802.11n/ac/ax device; a \$5 ESP32-S3 reference implementation produces 45--90~MB/month at zero marginal cost, compared to IBM~Quantum at \$1.60/second.
We provide a formal indistinguishability game and proof sketch for PUEK security under a spatial decorrelation assumption.
All code, data references, and the full extraction pipeline are open-source.
\end{abstract}

\begin{IEEEkeywords}
WiFi CSI, entropy source, NIST SP 800-90B, unilateral randomness, physical unclonable environment key, post-quantum cryptography, IoT security
\end{IEEEkeywords}

%% Body: include from the ACM version via input
%% All \section through \section*{Reproducibility} are format-agnostic
\input{body-ieee}

\section*{Acknowledgments}
The author thanks the TU~Darmstadt SEEMOO Lab and the University of Brescia for the public Gi-z/CSI-Data corpus, IBM Quantum for access to ibm\_kingston (156-qubit Heron~r2) via the IBM Quantum open plan, and Rigetti Computing for access to Ankaa-class processors used in entropy harvesting. Norwegian Patent Application filed April 2026 (Patentstyret) covers the CSI entropy extraction method and PUEK construction.

\bibliographystyle{IEEEtran}
\bibliography{references}

\end{document}
